\chaptercustom{4}{MÔ HÌNH GIẢI TÍCH VÀ KỊCH BẢN ĐO}
\setcounter{section}{4}
\setcounter{equation}{0}
Trước khi đo, nhóm xây dựng một mô hình giải tích đơn giản để dự đoán thời gian mỗi pha của vòng đồng bộ. Mô hình có hai tác dụng: cho biết trước kết quả ``đúng'' nên trông như thế nào để phát hiện lỗi cấu hình khi đo, và giúp giải thích phần chênh giữa số đo với lý thuyết. Chương này trình bày mô hình, các kịch bản đo và quy trình thu thập số liệu.

\subsection{Mô hình thời gian một pha truyền}
Xét $N$ luồng TCP song song, mỗi luồng truyền $S$ byte qua một link nghẽn chung có băng thông $B$, RTT cơ sở $R$. Giả thiết:
\begin{itemize}
    \item mỗi luồng bắt đầu slow start từ $IW = 10 \times 1448 = 14\,480$ byte (kết nối mới, hoặc sau khoảng nghỉ dài hơn RTO);
    \item tổng cửa sổ của $N$ luồng gấp đôi sau mỗi RTT cho tới khi đạt BDP, sau đó link được dùng hết ở tốc độ $\eta B$;
    \item không có mất gói ngoài lúc hàng đợi tràn, và các luồng chia đều băng thông.
\end{itemize}

Số RTT trong slow start và lượng dữ liệu đã gửi trong thời gian đó:
\begin{equation}
    k^{*} = \max\!\left(0,\ \left\lceil \log_2 \frac{B R / 8}{N \cdot IW} \right\rceil\right), \qquad
    D_{ss} = N \cdot IW \cdot \left(2^{k^{*}} - 1\right)
    \label{eq:ss}
\end{equation}

Nếu $D_{ss} < N S$, phần còn lại được truyền ở tốc độ đầy đủ:
\begin{equation}
    T_{phase}(N, S) = k^{*} R + \frac{8\,(N S - D_{ss})}{\eta B} + \frac{R}{2}
    \label{eq:tphase}
\end{equation}
Số hạng $R/2$ là thời gian để byte cuối cùng lan tới đích. Nếu khối đủ nhỏ để kết thúc ngay trong slow start ($D_{ss} \geq N S$), thời gian xấp xỉ $\lceil \log_2(S/IW + 1) \rceil R + R/2$.

Thời gian một vòng:
\begin{equation}
    T_{round} = T_{down} + T_{train} + T_{up} + T_{agg}
    \label{eq:tround}
\end{equation}
với $T_{down} = T_{phase}(N, S_{down})$, $T_{up} = T_{phase}(N, S_{up})$, $T_{train} = 1$~s và $T_{agg} \approx 0$.

Khi $N S$ lớn, số hạng thứ hai của~\eqref{eq:tphase} chiếm ưu thế và
\begin{equation}
    T_{phase} \approx \frac{8 N S}{\eta B}
    \label{eq:linear}
\end{equation}
tăng tuyến tính theo $N$ với hệ số góc $8S/(\eta B)$. Với $S = 10$~MB và $B = 100$~Mbps, hệ số góc là $0{,}837$~s cho mỗi Edge Node thêm vào, \emph{mỗi pha}.

Mô hình được cài đặt trong \texttt{code/model.py}; kết quả nằm trong thư mục \texttt{data/} dưới dạng các file \texttt{model\_<vị trí>\_<kích thước>.csv} và được dùng trực tiếp để vẽ biểu đồ ở Chương~5.

Mô hình không tính hành vi răng cưa của CUBIC sau mất gói, độ trễ hàng đợi làm RTT thực tế lớn hơn $R$, chi phí CPU của Mininet, và việc các luồng không bắt đầu đồng thời hoàn toàn. Vì vậy số đo được kỳ vọng lớn hơn mô hình một chút; nếu số đo nhỏ hơn mô hình đáng kể, nhiều khả năng cấu hình băng thông chưa có hiệu lực.

\subsection{Kịch bản đo}
Năm kịch bản trong Bảng~\ref{tab:scen} được cài sẵn trong \texttt{code/run\_experiments.py}; mỗi kịch bản chạy với $N = 2, 3, 4, 5, 6$.

\begin{table}[H]
    \centering
    \caption{Các kịch bản đo}
    \label{tab:scen}
    \begin{tabular}{l l c c c l}
        \toprule
        \textbf{Mã} & \textbf{Aggregator} & \textbf{WAN} & $S_{down}$ & $S_{up}$ & \textbf{Phục vụ} \\
        \midrule
        K1\_edge\_full & Edge & -- & 10 MB & 10 MB & Q1, Q2, Q3 \\
        K2\_edge\_comp & Edge & -- & 10 MB & 1 MB & Q2 \\
        K2\_cloud\_comp & Cloud & 100 Mbps & 10 MB & 1 MB & Q2 \\
        K3\_cloud\_full & Cloud & 100 Mbps & 10 MB & 10 MB & Q2, Q3 \\
        K3\_cloud20\_full & Cloud & 20 Mbps & 10 MB & 10 MB & Q3 \\
        \bottomrule
    \end{tabular}
\end{table}

Ở kịch bản nén, chỉ chiều lên được nén vì đây là cách các kỹ thuật nén gradient thường áp dụng: máy khách gửi gradient thưa, còn mô hình toàn cục phát xuống vẫn đầy đủ.

\subsection{Quy trình đo}
\begin{enumerate}
    \item Khởi tạo topo mới cho mỗi lần đo, tắt TSO/GSO/GRO trên mọi giao diện.
    \item Bật \texttt{ping -i 0.2} từ \texttt{h1} tới nút tổng hợp trong suốt lần đo để ghi RTT khi có tải.
    \item (Tuỳ chọn \texttt{-{}-pcap}) bật \texttt{tcpdump} trên giao diện của nút tổng hợp.
    \item Chạy \texttt{fl\_server.py} trên nút tổng hợp, sau đó $N$ tiến trình \texttt{fl\_client.py} trên \texttt{h1..hN}.
    \item Server ghi mỗi vòng một dòng gồm $T_{down}$, $T_{up}$, $T_{round}$ vào \texttt{results/<kịch bản>\_n<N>.csv}.
    \item Sau 10 vòng, dừng mọi tiến trình và huỷ topo.
    \item Chạy \texttt{code/summarize.py}: bỏ vòng 1 (ARP, bắt tay, khởi động Python), lấy trung bình và độ lệch chuẩn của 9 vòng còn lại, ghi vào \texttt{data/measured\_<kịch bản>.csv}.
\end{enumerate}

Lệnh chạy toàn bộ:
\codesnap{run_all}

\subsection{Cách đo từng chỉ số}
\begin{itemize}
    \item $T_{down}$: từ lúc server bắt đầu gửi \texttt{GLOBAL} tới lúc nhận \texttt{ACK} tầng ứng dụng của client cuối cùng.
    \item $T_{up}$: từ lúc client đầu tiên kết thúc huấn luyện tới lúc \texttt{UPDATE} cuối cùng tới server, tính bằng $\max(t_{update}) - \min(t_{ack}) - T_{train}$.
    \item $T_{round}$: từ lúc bắt đầu gửi \texttt{GLOBAL} tới lúc server gửi xong \texttt{ACK} kết thúc vòng.
    \item RTT khi có tải: giá trị \texttt{ping} trong khoảng thời gian có truyền dữ liệu, so với RTT lúc rảnh.
    \item $d_{trans}$ hiệu dụng: $T_{up}$ trừ phần slow start đọc từ biểu đồ Stevens.
\end{itemize}

\subsection{Các yếu tố ảnh hưởng độ tin cậy}
Mininet chạy mọi host trên cùng một nhân hệ điều hành, nên CPU của máy ảo là nguồn sai số lớn nhất. Nhóm cấp ít nhất 2 vCPU và chỉ đo khi \texttt{iperf} đạt gần 95 Mbps. Đồng hồ thì không gây sai số: mọi thời điểm đều lấy trên server.

Hai nguồn sai số còn lại nằm ở cách đo. Nếu lấy $T_{down}$ theo thời điểm \texttt{sendall()} trả về, kết quả sẽ nhỏ hơn thực tế vài trăm ms; ACK tầng ứng dụng của FLSync có để tránh lỗi này. Kích thước hàng đợi (\nolinkurl{max_queue_size}) cũng ảnh hưởng: quá lớn thì RTT khi có tải tăng mạnh, quá nhỏ thì mất gói liên tục. Nhóm chọn 200 gói, tức tối đa khoảng 24~ms trễ hàng đợi ở 100 Mbps.

\subsection{Kết luận chương}
Mô hình giải tích dự đoán thời gian mỗi pha tăng tuyến tính theo $N$ với hệ số góc khoảng 0{,}84~s mỗi Edge Node cho khối 10 MB trên link 100 Mbps, cộng thêm một phần cố định do slow start phụ thuộc RTT. Năm kịch bản đo được thiết kế để kiểm chứng từng câu hỏi Q1--Q3 với cùng một quy trình tự động.

\cleardoublepage
